\documentclass[11pt]{article}
\usepackage[a4paper,margin=1in]{geometry}
\usepackage{fontspec}
\usepackage[boldfont=false]{xeCJK}
\setCJKmainfont{FandolSong-Regular.otf}
\usepackage{amsmath}
\usepackage{amssymb}
\usepackage{array}
\usepackage{booktabs}
\usepackage{graphicx}
\usepackage{adjustbox}
\usepackage{enumitem}
\setlist[itemize]{topsep=2pt,partopsep=0pt,parsep=0pt,itemsep=2pt,leftmargin=1.4em}
\setlist[enumerate]{topsep=2pt,partopsep=0pt,parsep=0pt,itemsep=2pt,leftmargin=1.6em}
\usepackage{parskip}
\usepackage{fvextra}
\fvset{breaklines=true,breakanywhere=true,fontsize=\small}
\setlength{\parindent}{0pt}

\begin{document}

\begin{center}
  {\LARGE\bfseries Organic synthesis}\\[0.35em]
  {\large A Level Chemistry}
\end{center}
\vspace{0.6em}

\section*{Organic synthesis}

Like the AS synthesis topic, this asks you to \textbf{join up} known reactions to build a target molecule. Now you also have the A Level reactions of arenes, phenol, amines, amides and acyl chlorides.

\subsection*{Identifying functional groups}

A molecule may carry several \textbf{functional group} 官能团 types. Use the test reactions to spot each one, then predict its behaviour. For example:

\begin{itemize}
\item decolourises bromine water with \textbf{no} catalyst, giving a white precipitate → a phenol or a phenylamine (the ring is activated).

\item reacts violently with cold water, giving fumes of $\text{HCl}$ → an \textbf{acyl chloride} 酰氯.

\item gives a purple colour with neutral $\text{FeCl}_3$ → a \textbf{phenol} 苯酚.

\end{itemize}

\subsection*{A map of the A Level reactions}

\begin{center}
\adjustbox{max width=\linewidth}{%
\begin{tabular}{lll}
\toprule
\textbf{Start} & \textbf{Reagent and conditions} & \textbf{Product} \\
\midrule
benzene (an \textbf{arene} 芳烃) & $\text{HNO}_3$ / $\text{H}_2\text{SO}_4$ (\textbf{electrophilic substitution} 亲电取代) & nitrobenzene \\
nitrobenzene & $\text{Sn}$ / conc $\text{HCl}$, then $\text{NaOH}$ (\textbf{reduction} 还原) & phenylamine \\
phenylamine & $\text{HNO}_2$, below $10\,°\text{C}$ & a diazonium salt \\
diazonium salt & warm water & phenol; or couple with phenol → azo dye \\
methylbenzene & hot $\text{KMnO}_4$ (\textbf{oxidation} 氧化) & benzoic acid \\
\textbf{carboxylic acid} 羧酸 & $\text{SOCl}_2$ or $\text{PCl}_5$ & acyl chloride \\
acyl chloride & alcohol / phenol & an \textbf{ester} 酯 \\
acyl chloride & ammonia / \textbf{amine} 胺 & an \textbf{amide} 酰胺 \\
amide or \textbf{nitrile} 腈 & $\text{LiAlH}_4$ & an amine \\
halogenoalkane & $\text{KCN}$ & a nitrile (adds one carbon) \\
\bottomrule
\end{tabular}}
\end{center}

\subsection*{Planning and analysing a route}

To devise a \textbf{synthetic route} 合成路线, work \textbf{backwards} from the target: which single reaction makes it, and from what? Repeat until you reach the starting material, then write each step with its \textbf{reagent} 试剂 and conditions.

When you \textbf{analyse} a route, state the type of reaction for each step (for example electrophilic substitution, \textbf{addition–elimination} 加成消去, oxidation or reduction) and watch for likely \textbf{by-products} 副产物 — for example, making an amine from a halogenoalkane also gives over-substituted amines, lowering the yield.


\end{document}
